% Корректировка горизонтального положения числителя и знаменателя
\directlua{
    left_spacings = {
        A = 0.05, H = 0.05, N = 0.025, a = 0.05, b = 0.05, c = 0.05,
        e = 0.025, f = 0.15, g = 0.1, h = 0.05, i = 0.05, j = 0.2,
        m = 0.05, n = 0.05, o = 0.05, p = 0.1, q = 0.05, r = 0.025,
        s = 0.075, t = 0.05, x = 0.05, z = 0.05, y = 0.075, z = 0.05
    }
    right_spacings = {
        C = -0.045, E = -0.031, F = -0.065, G = -0.012,
        H = -0.095, I = -0.085, J = -0.11, K = -0.07,
        M = -0.1, N = -0.1, O = -0.005, P = -0.02,
        S = -0.05, T = -0.1, U = -0.1,
        V = -0.12, W = -0.12, X = -0.135, Y = -0.12,
        Z = -0.06, d = -0.05, f = -0.05, g = -0.03,
        j = -0.01, k = -0.05, l = -0.025, r = -0.025,
        t = -0.02, v = -0.005, z = -0.01
    }

    function match_first_letter(math_token)
        return math_token:match("^ *\\subpow *[^ a-zA-Z] *([a-zA-Z])")
            or math_token:match("^ *\\pow *[^ a-zA-Z] *([a-zA-Z])")
            or math_token:match("^ *\\sub *[^ a-zA-Z] *([a-zA-Z])")
            or math_token:match("^ *([a-zA-Z])")
            or ""
        end

    function match_last_letter(math_token)
        return math_token:match("([a-zA-Z]) *$") or ""
    end

    function compute_dfrac(indent, numerator, denominator, vphantom)
        indent = "\\hspace{" .. indent .. "}"
        local numerator_left_spacing = left_spacings[match_first_letter(numerator)] or 0
        local numerator_right_spacing = right_spacings[match_last_letter(numerator)] or 0
        local denominator_left_spacing = left_spacings[match_first_letter(denominator)] or 0
        local denominator_right_spacing = right_spacings[match_last_letter(denominator)] or 0
        if numerator_left_spacing == 0 then else
            numerator = "\\kern" .. numerator_left_spacing .. "em " .. numerator
        end
        if numerator_right_spacing == 0 then else
            numerator = numerator .. "\\kern" .. numerator_right_spacing .. "em "
        end
        if denominator_left_spacing >= 0 then else
            denominator = "\\kern" .. denominator_left_spacing .. "em " .. denominator
        end
        if denominator_right_spacing >= 0 then else
            denominator = denominator .. "\\kern" .. denominator_right_spacing .. "em "
        end
        if vphantom == "" then else
            numerator = numerator .. "\\vphantom{" .. vphantom .. "}"
        end
        tex.print("\\olddfrac{" .. indent .. numerator .. indent .. "}{" .. indent .. denominator .. indent .. "}")
    end
}

\AtBeginDocument{
    % Запрет перехода на следующую страницу внутри абзаца
    \widowpenalties 1 10000
    \raggedbottom

    % Запрет переноса внутри математического окружения
    \binoppenalty=10000
    \relpenalty=10000

    % Принудительный отступ между строками
    \setlength{\lineskiplimit}{1mm}
    \setlength{\lineskip}{1mm}
    % Установить отступы перед и после математических окружений (align, multline, gather и т.д.)
    \setlength{\abovedisplayskip}{0.5mm}
    \setlength{\belowdisplayskip}{0.5mm}
    \setlength{\abovedisplayshortskip}{0.25mm}
    \setlength{\belowdisplayshortskip}{0.25mm}

    % Символы с фиксированными отступами

    % Знак прим и его вариант для строчных букв
    \renewmathcommand{\prime}{\char"2032\kern-0.06em}
    \newmathcommand{\lprime}{\raisebox{-0.15em}{$\prime$}}
    % Сокращения для предыдущих, которые могут не работать внутри
    % некоторых окружений, например в framed.
    \renewmathcommand{\'}{\prime}
    \renewmathcommand{\`}{\lprime}
    % Знак прим с нижним индексом
    \newmathcommand{\primesub}[3][1]{
        \mathrlap{\hphantom{#2}\multido{}{#1}{\prime}\kern-0.06em}
        \subpow{#2}{#3}{\hphantom{\displaystyle\multido{}{#1}{\prime}}\kern-0.06em}
    }

    \newmathcommand{\f}{f\hspace{0.075em}}
    \renewmathcommand{\j}{\hspace{0.15em}j}
    \newmathcommand{\x}{\hspace{0.05em} x}
    \newmathcommand{\y}{\hspace{0.05em} y}
    \newmathcommand{\p}{\ \! p}
    \newmathcommand{\dfracphantom}{\vphantom{y}}

    \let\olddfrac\dfrac
    \renewcommand{\dfrac}[3][0.1em]{
        \directlua{
            local indent = "\luaescapestring{\unexpanded{#1}}"
            local numerator = "\luaescapestring{\unexpanded{#2}}"
            local denominator = "\luaescapestring{\unexpanded{#3}}"
            compute_dfrac(indent, numerator, denominator, "")
        }
    }

    % Дробь с вертикальным отсупом (вставляется вертикальный фантом с буквой y)
    \newmathcommand{\dvfrac}[3][0.1em]{
        \directlua{
            local indent = "\luaescapestring{\unexpanded{#1}}"
            local numerator = "\luaescapestring{\unexpanded{#2}}"
            local denominator = "\luaescapestring{\unexpanded{#3}}"
            compute_dfrac(indent, numerator, denominator, "y")
        }
    }

    \let\oldfrac\frac
    \renewcommand{\frac}[3][0.1em]{\oldfrac{\kern#1 #2 \kern #1}{\kern#1 #3 \kern#1}}
    \newcommand{\vfrac}[3][0.1em]{\oldfrac{\kern#1 #2 \vphantom{y} \kern #1}{\kern#1 #3 \kern#1}}
}